% !TeX program = pdflatex
\documentclass[10pt, aspectratio=169]{beamer}

% Preámbulo modular para presentaciones Beamer (Modelación Estadística)
% Diseño y paleta institucional del Tecnológico de Monterrey

\documentclass[aspectratio=169,xcolor={dvipsnames,table}]{beamer}

\usepackage[utf8]{inputenc}
\usepackage[spanish,mexico]{babel}
\usepackage{amsmath,amssymb,amsfonts,mathtools}
\usepackage{graphicx}
\usepackage{listings}
\usepackage{booktabs}
\usepackage{tikz}
\usetikzlibrary{matrix,arrows,decorations.pathmorphing,shapes.geometric,positioning}

% Paleta Institucional Tecnológico de Monterrey
\definecolor{TecAzul}{HTML}{0039A6}       % Azul TEC: color primario institucional
\definecolor{TecAzulOscuro}{HTML}{1A2E51} % Azul marino oscuro
\definecolor{TecRojo}{HTML}{EC2661}       % Rojo de acento (paleta miTec)
\definecolor{TecGris}{HTML}{646464}       % Gris neutro para comentarios y subtítulos
\definecolor{TecFondoBloque}{HTML}{F4F6F9}% Fondo suave para bloques

% Configuración de Tema Beamer
\usetheme{Boadilla}
\usecolortheme[named=TecAzulOscuro]{structure}
\setbeamercolor{alerted text}{fg=TecRojo}
\setbeamercolor{palette primary}{bg=TecAzulOscuro,fg=white}
\setbeamercolor{palette secondary}{bg=TecAzul,fg=white}
\setbeamercolor{palette tertiary}{bg=TecAzulOscuro!80!black,fg=white}
\setbeamercolor{title}{fg=TecAzulOscuro,bg=TecFondoBloque}
\setbeamercolor{frametitle}{fg=TecAzulOscuro,bg=TecFondoBloque}

% Bloques personalizados elegantes
\setbeamercolor{block title}{bg=TecAzulOscuro,fg=white}
\setbeamercolor{block body}{bg=TecFondoBloque,fg=black}
\setbeamercolor{block title alerted}{bg=TecRojo,fg=white}
\setbeamercolor{block body alerted}{bg=TecRojo!8,fg=black}
\setbeamercolor{block title example}{bg=TecAzul,fg=white}
\setbeamercolor{block body example}{bg=TecAzul!8,fg=black}

% Redondear bordes y añadir sombra sutil
\setbeamertemplate{blocks}[rounded][shadow=true]
\setbeamertemplate{navigation symbols}{}

% Pie de página limpio con número de diapositiva
\setbeamertemplate{footline}{
  \leavevmode%
  \hbox{%
  \begin{beamercolorbox}[wd=.7\paperwidth,ht=2.25ex,dp=1ex,left]{author in head/foot}%
    \hspace*{2ex}\insertshortauthor~~(\insertshortinstitute) \hfill \insertshorttitle
  \end{beamercolorbox}%
  \begin{beamercolorbox}[wd=.3\paperwidth,ht=2.25ex,dp=1ex,right]{date in head/foot}%
    \insertshortdate{}\hspace*{2em}
    \insertframenumber{} / \inserttotalframenumber\hspace*{2ex}
  \end{beamercolorbox}}%
  \vskip0pt%
}

% Configuración de Listings de Python para visualización pedagógica de código
\lstset{
    language=Python,
    basicstyle=\ttfamily\footnotesize,
    keywordstyle=\color{TecAzulOscuro}\bfseries,
    stringstyle=\color{TecRojo},
    commentstyle=\color{TecGris}\itshape,
    showstringspaces=false,
    breaklines=true,
    frame=lines,
    rulecolor=\color{TecAzulOscuro!30},
    backgroundcolor=\color{TecFondoBloque},
    aboveskip=0.5em,
    belowskip=0.5em,
    literate={á}{{\'a}}1 {é}{{\'e}}1 {í}{{\'i}}1 {ó}{{\'o}}1 {ú}{{\'u}}1
             {Á}{{\'A}}1 {É}{{\'E}}1 {Í}{{\'I}}1 {Ó}{{\'O}}1 {Ú}{{\'U}}1
             {ñ}{{\~n}}1 {Ñ}{{\~N}}1 {¿}{{?`}}1 {¡}{{!`}}1
}

% Metadatos institucionales por defecto
\title[Modelación Estadística]{Modelación Estadística para Ciencia de Datos e Ingeniería}
\author[J. Castillo Colmenares]{Juliho David Castillo Colmenares}
\institute[Tec de Monterrey]{Tecnológico de Monterrey \\ Escuela de Ingeniería y Ciencias}
\date{\today}

% Comandos y macros estadísticas compartidas para las presentaciones Beamer (Metropolis)

% Conjuntos numéricos básicos
\newcommand{\R}{\mathbb{R}}
\newcommand{\N}{\mathbb{N}}
\newcommand{\Q}{\mathbb{Q}}
\newcommand{\Z}{\mathbb{Z}}
\newcommand{\set}[1]{\left\{ #1 \right\}}
\newcommand{\abs}[1]{\left|#1\right|}
\newcommand{\norm}[1]{\left\| #1 \right\|}

% Comandos estadísticos y probabilísticos del curso
\newcommand{\Var}{\operatorname{Var}}
\newcommand{\cov}{\operatorname{Cov}}
\newcommand{\corr}{\rho}
\newcommand{\comb}[2]{\begin{pmatrix} #1 \\ #2 \end{pmatrix}}
\newcommand{\s}{\sigma}
\newcommand{\particion}{\mathfrak{P}}
\newcommand{\card}[1]{\#\left( #1 \right)}
\newcommand{\seq}[3]{\set{#1}_{#2}^{#3}}
\newcommand{\E}{\mathbb{E}}
\newcommand{\Prob}{\mathbb{P}}

% Entornos pedagógicos y cajas en español académico natural
\newenvironment{discusion}{%
  \begin{alertblock}{Pregunta de Discusión}%
}{%
  \end{alertblock}%
}

\newenvironment{reflexion}{%
  \begin{alertblock}{Reflexión para el Aula}%
}{%
  \end{alertblock}%
}

\newenvironment{paradoja}{%
  \begin{alertblock}{Paradoja de Exploración}%
}{%
  \end{alertblock}%
}

\newenvironment{motivacion}{%
  \begin{exampleblock}{Motivación: Ciencia de Datos en la Práctica}%
}{%
  \end{exampleblock}%
}

\newenvironment{retoclase}{%
  \begin{block}{Reto Rápido en Equipo}%
}{%
  \end{block}%
}


\title[Esperanza y Varianza]{Sección 03.03: Esperanza y Varianza}
\subtitle{Momentos Operacionales y Centro de Gravedad Probabilístico}
\author[J. Castillo Colmenares]{Juliho Castillo Colmenares}
\institute{Tecnológico de Monterrey}
\date{\vspace{-2.0cm}}

\begin{document}

% Slide 1: Portada
\begin{frame}[plain]
  \vspace{-0.5cm}
  \titlepage
\end{frame}

% Slide 2: Hoja de Ruta
\begin{frame}{Hoja de Ruta — Capítulo 03: Variables Aleatorias Discretas}
  \begin{block}{Estructura Curricular del Capítulo 03}
    \footnotesize
    \begin{enumerate}\itemsep=0.08cm
      \item \textcolor{gray}{Sección 03.01: Funciones de Probabilidad Discretas (PMF y Soporte)}
      \item \textcolor{gray}{Sección 03.02: Función de Distribución Acumulada (CDF Discreta)}
      \item \textbf{\textcolor{TecRojo}{Sección 03.03: Esperanza Matemática, Varianza y Momentos}} $\leftarrow$ \emph{Foco actual}
      \item \textcolor{gray}{Sección 03.04: Distribuciones de Bernoulli y Binomial}
      \item \textcolor{gray}{Sección 03.05: Distribuciones Geométrica y Binomial Negativa}
      \item \textcolor{gray}{Sección 03.06: Distribución Hipergeométrica}
    \end{enumerate}
  \end{block}
  \vspace{-0.05cm}
  \pause
  \begin{alertblock}{Objetivo Curricular de la Sección}
    \footnotesize
    Dominar la fundamentación algebraica y propiedades operacionales del valor esperado y la varianza, aplicando el Teorema LOTUS y el teorema del centro de gravedad probabilístico en problemas de toma de decisiones e ingeniería.
  \end{alertblock}
\end{frame}

% Slide 3: Motivación e Intuición
\begin{frame}{Motivación — De la Distribución Completa a los Resúmenes Numéricos}
  \begin{columns}[T]
    \begin{column}{0.48\textwidth}
      \begin{block}{La Complejidad de la Distribución Exhaustiva}
        \footnotesize
        \begin{itemize}\itemsep=0.1cm
          \item Una función de masa $f_X(x)$ o acumulada $F_X(x)$ contiene la \textbf{información estocástica completa} del fenómeno.
          \item En ingeniería de sistemas y finanzas, comparar distribuciones completas resulta computacional o algebraicamente abrumador.
          \item Se requieren \textbf{escalares sintéticos invariantes} que resuman la posición central y la dispersión del riesgo.
        \end{itemize}
      \end{block}
    \end{column}
    \begin{column}{0.48\textwidth}
      \pause
      \begin{alertblock}{Centro de Masa e Inercia Físico-Estocástica}
        \footnotesize
        \begin{itemize}\itemsep=0.1cm
          \item \textbf{Esperanza Matemática ($\mu$):} Es el punto de equilibrio mecánico o centro de masa poblacional de la distribución. Pondera cada valor por su probabilidad.
          \item \textbf{Varianza ($\sigma^2$):} Es el momento de inercia respecto a $\mu$. Cuantifica el riesgo, la volatilidad y la dispersión cuadrática del sistema.
        \end{itemize}
      \end{alertblock}
    \end{column}
  \end{columns}
\end{frame}

% Slide 4: Definición Formal de Esperanza Matemática
\begin{frame}{Definición Formal — Esperanza Matemática}
  \begin{block}{Definición Analítica del Valor Esperado ($\mu = \E[X]$)}
    \footnotesize
    Sea $X$ una variable aleatoria discreta con soporte numerable $\mathcal{S}_X = \{x_1, x_2, \dots\}$ y función de masa de probabilidad $f_X(x)$. La \textbf{esperanza matemática} o valor esperado de $X$ se define como:
    $$\mu = \E[X] = \sum_{x \in \mathcal{S}_X} x \cdot f_X(x),$$
    siempre que la serie converga absolutamente ($\sum |x| f_X(x) < \infty$).
  \end{block}
  \vspace{-0.05cm}
  \pause
  \begin{columns}[T]
    \begin{column}{0.48\textwidth}
      \begin{exampleblock}{Interpretación por la Ley de Grandes Números}
        \footnotesize
        Si un experimento se repite de forma independiente $n \to \infty$ veces, el promedio aritmético muestral $\bar{X}_n$ converge casi seguramente a la media poblacional $\E[X]$.
      \end{exampleblock}
    \end{column}
    \begin{column}{0.48\textwidth}
      \begin{alertblock}{Advertencia sobre el Soporte}
        \footnotesize
        El valor $\E[X]$ \textbf{no tiene por qué pertenecer} al soporte $\mathcal{S}_X$. Por ejemplo, en el lanzamiento de un dado balanceado, $\E[X] = 3.5 \notin \{1, 2, \dots, 6\}$.
      \end{alertblock}
    \end{column}
  \end{columns}
\end{frame}

% Slide 5: Propiedades Operacionales y Teorema LOTUS
\begin{frame}{Propiedades Operacionales y Teorema LOTUS}
  \begin{block}{Teorema del Estadístico Inconsciente (LOTUS)}
    \footnotesize
    Sea $g: \mathbb{R} \to \mathbb{R}$ una función real y $Y = g(X)$ una variable aleatoria transformada. El valor esperado de $Y$ se puede obtener directamente sobre el soporte original de $X$:
    $$\E[g(X)] = \sum_{x \in \mathcal{S}_X} g(x) \cdot f_X(x),$$
    sin necesidad de derivar previamente la distribución inducida $f_Y(y)$.
  \end{block}
  \vspace{-0.05cm}
  \pause
  \begin{columns}[T]
    \begin{column}{0.48\textwidth}
      \begin{exampleblock}{Linealidad Universal del Operador Esperanza}
        \footnotesize
        Para cualesquiera constantes $c, d \in \mathbb{R}$ y variables $X, Y$:
        $$\E[cX + dY] = c\E[X] + d\E[Y]$$
        Esta propiedad es válida \textbf{incluso si las variables son dependientes}.
      \end{exampleblock}
    \end{column}
    \begin{column}{0.48\textwidth}
      \begin{alertblock}{Monotonía y Cota de Esperanza}
        \footnotesize
        Si $X \ge 0$ casi seguramente, entonces $\E[X] \ge 0$. Si $a \le X \le b$ en todo su soporte, entonces $a \le \E[X] \le b$.
      \end{alertblock}
    \end{column}
  \end{columns}
\end{frame}

% Slide 6: Varianza Operacional y Fórmula de König-Huygens
\begin{frame}{Varianza Operacional y Fórmula de König-Huygens}
  \begin{block}{Definición Central de la Varianza ($\sigma^2 = \Var(X)$)}
    \footnotesize
    La varianza es el segundo momento central con respecto a la esperanza $\mu = \E[X]$:
    $$\sigma^2 = \Var(X) = \E\left[ (X - \mu)^2 \right] = \sum_{x \in \mathcal{S}_X} (x - \mu)^2 \cdot f_X(x) \ge 0.$$
    Su raíz cuadrada positiva $\sigma = \sqrt{\Var(X)}$ se denomina \textbf{desviación estándar} y se expresa en las mismas unidades físicas que $X$.
  \end{block}
  \vspace{-0.05cm}
  \pause
  \begin{columns}[T]
    \begin{column}{0.48\textwidth}
      \begin{alertblock}{Fórmula de König-Huygens (Atajo Operacional)}
        \footnotesize
        Expandiendo el cuadrado y aplicando linealidad:
        $$\Var(X) = \E[X^2 - 2\mu X + \mu^2] = \E[X^2] - \mu^2$$
        Es el método preferido en cálculo manual o algorítmico.
      \end{alertblock}
    \end{column}
    \begin{column}{0.48\textwidth}
      \begin{block}{Propiedades de Escala e Independencia}
        \footnotesize
        $$\Var(cX + d) = c^2 \Var(X)$$
        Si $X, Y$ son independientes: $\Var(X \pm Y) = \Var(X) + \Var(Y)$.
      \end{block}
    \end{column}
  \end{columns}
\end{frame}

% Slide 7: Taxonomía de Momentos y Asimetría
\begin{frame}{Taxonomía de Momentos, Asimetría y Curtosis}
  \begin{columns}[T]
    \begin{column}{0.48\textwidth}
      \begin{block}{Momentos Brutos vs Momentos Centrales}
        \footnotesize
        \begin{itemize}\itemsep=0.1cm
          \item \textbf{Momento Bruto de orden $k$ ($\mu_k'$):}
          $$\mu_k' = \E[X^k] = \sum_{x \in \mathcal{S}_X} x^k f_X(x)$$
          \item \textbf{Momento Central de orden $k$ ($\mu_k$):}
          $$\mu_k = \E\left[(X - \mu)^k\right] = \sum_{x \in \mathcal{S}_X} (x - \mu)^k f_X(x)$$
        \end{itemize}
      \end{block}
    \end{column}
    \begin{column}{0.48\textwidth}
      \pause
      \begin{alertblock}{Asimetría ($\gamma_1$) y Curtosis Excesiva ($\gamma_2$)}
        \footnotesize
        $$\gamma_1 = \dfrac{\mu_3}{\sigma^3} = \dfrac{\E\left[(X-\mu)^3\right]}{\sigma^3}$$
        $$\gamma_2 = \dfrac{\mu_4}{\sigma^4} - 3 = \dfrac{\E\left[(X-\mu)^4\right]}{\sigma^4} - 3$$
        $\gamma_1$ mide el sesgo lateral y $\gamma_2$ el peso relativo en las colas respecto a una distribución normal.
      \end{alertblock}
    \end{column}
  \end{columns}
\end{frame}

% Slide 8: Estandarización y Centro de Gravedad Probabilístico
\begin{frame}{Estandarización ($Z$-Scores) y Centro de Gravedad Probabilístico}
  \begin{columns}[T]
    \begin{column}{0.48\textwidth}
      \begin{block}{Variable Aleatoria Estandarizada ($Z$-Score)}
        \footnotesize
        Para cualquier variable discreta $X$ con media $\mu$ y varianza finita $\sigma^2 > 0$, su transformación estandarizada es:
        $$Z = \dfrac{X - \mu}{\sigma}$$
        Cumple universalmente que $\E[Z] = 0$ y $\Var(Z) = 1$, permitiendo comparar escalas heterogéneas.
      \end{block}
    \end{column}
    \begin{column}{0.48\textwidth}
      \pause
      \begin{alertblock}{Teorema del Centro de Gravedad}
        \footnotesize
        El valor esperado $\mu = \E[X]$ es el único escalar real que minimiza el error cuadrático medio de aproximación:
        $$\min_{a \in \mathbb{R}} \E\left[ (X - a)^2 \right] = \Var(X),$$
        alcanzándose el mínimo estrictamente en $a^* = \mu$.
      \end{alertblock}
    \end{column}
  \end{columns}
\end{frame}

% Slide 9A: Laboratorio en Python — Momentos
\begin{frame}[fragile]{Laboratorio en Python: Cálculo Exacto de Momentos (`numpy`)}
  \lstinputlisting[language=Python, firstline=12, lastline=28, basicstyle=\fontsize{6.5pt}{7.5pt}\selectfont\ttfamily]{../../code/03_variables_aleatorias_discretas/03.03_expectation_and_variance.py}
\end{frame}

% Slide 9B: Laboratorio en Python — LOTUS & Z-Score
\begin{frame}[fragile]{Laboratorio en Python: Verificación LOTUS y Estandarización $Z$}
  \lstinputlisting[language=Python, firstline=31, lastline=50, basicstyle=\fontsize{6.5pt}{7.5pt}\selectfont\ttfamily]{../../code/03_variables_aleatorias_discretas/03.03_expectation_and_variance.py}
\end{frame}

% Slide 9C: Laboratorio en Python — Monte Carlo
\begin{frame}[fragile]{Laboratorio en Python: Simulación Monte Carlo de la Ley de Grandes Números}
  \lstinputlisting[language=Python, firstline=53, lastline=67, basicstyle=\fontsize{6.5pt}{7.5pt}\selectfont\ttfamily]{../../code/03_variables_aleatorias_discretas/03.03_expectation_and_variance.py}
\end{frame}

% Slide 9D: Laboratorio en Python — Salida en Consola
\begin{frame}[fragile]{Laboratorio en Python: Salida en Consola}
  \begin{block}{Ejecución en Consola (\texttt{python 03.03\_expectation\_and\_variance.py})}
    \fontsize{6.5pt}{7.5pt}\selectfont\ttfamily
=== Section 03.03: Expectation and Variance Lab ===\\[0.1cm]
{[1]} Theoretical Moments:\\
\ \ \ \ Expectation (Mean) E[X]\ \ : 1.2000\\
\ \ \ \ Second Raw Moment\ \ E[X\^{}2]: 2.8000\\
\ \ \ \ Variance\ \ \ \ \ \ \ \ \ \ \ Var(X): 1.3600\\
\ \ \ \ Standard Deviation sigma : 1.1662\\[0.1cm]
{[2]} LOTUS \& Standardization Invariance:\\
\ \ \ \ Expected g(X) via LOTUS\ \ : 7.0000\\
\ \ \ \ Standardized Z Mean E[Z] : 0.000000 (Target: 0.0)\\
\ \ \ \ Standardized Z Var Var(Z): 1.000000 (Target: 1.0)\\[0.1cm]
{[3]} Monte Carlo LLN Simulation (N=100,000):\\
\ \ \ \ Empirical Sample Mean\ \ \ \ : 1.1955 (Error: 0.004460)\\
\ \ \ \ Empirical Sample Variance: 1.3532 (Error: 0.006796)\\[0.1cm]
=== All numerical verifications passed exact tolerances. ===
  \end{block}
\end{frame}

% Slide 10A: Ejercicio en Clase — Nivel 1: Fundamental (Enunciado)
\begin{frame}{Ejercicio en Clase — Nivel 1: Fundamental (1/2)}
  \begin{block}{Problema 3.3.1 — Transacciones en Error durante Hora Pico en BD Bancaria}
    \footnotesize
    Sea $X$ una variable aleatoria discreta que modela el número de transacciones con error durante una hora pico en un servidor bancario, con soporte $\mathcal{S}_X = \{0, 1, 2, 3, 4\}$ y función de masa explícita:
    \begin{table}[h]
      \centering
      \begin{tabular}{c|ccccc}
        \toprule
        $x$ & $0$ & $1$ & $2$ & $3$ & $4$ \\
        \midrule
        $f_X(x)$ & $0.35$ & $0.30$ & $0.20$ & $0.10$ & $0.05$ \\
        \bottomrule
      \end{tabular}
    \end{table}
    \begin{enumerate}\itemsep=0.1cm
      \item Calcule formalmente la esperanza matemática $\mu = \E[X]$ utilizando la definición analítica fundamental.
      \item Calcule el segundo momento no central $\E[X^2]$ y determine la varianza operacional $\Var(X)$ mediante la identidad de König-Huygens.
    \end{enumerate}
  \end{block}
\end{frame}

% Slide 10B: Ejercicio en Clase — Nivel 1: Fundamental (Resolución)
\begin{frame}{Ejercicio en Clase — Nivel 1: Fundamental (2/2)}
  \fontsize{6.8pt}{8.2pt}\selectfont
  \begin{block}{1. Esperanza Matemática ($\mu = \E[X]$)}
    \fontsize{6.8pt}{8.2pt}\selectfont
    Ponderamos cada valor de la variable por su probabilidad poblacional:
    $$\mu = \sum_{x=0}^4 x \cdot f_X(x) = 0(0.35) + 1(0.30) + 2(0.20) + 3(0.10) + 4(0.05)$$
    $$\mu = 0 + 0.30 + 0.40 + 0.30 + 0.20 = \mathbf{1.20\text{ transacciones}}$$
  \end{block}

  \vspace{-0.15cm}
  \pause
  \begin{alertblock}{2. Segundo Momento Bruto ($\E[X^2]$) y Varianza Operacional ($\Var(X)$)}
    \fontsize{6.8pt}{8.2pt}\selectfont
    Calculamos el promedio ponderado de los cuadrados del soporte:
    $$\E[X^2] = 0^2(0.35) + 1^2(0.30) + 2^2(0.20) + 3^2(0.10) + 4^2(0.05) = \mathbf{2.80}$$
    Aplicamos el atajo algebraico de König-Huygens:
    $$\Var(X) = \E[X^2] - \mu^2 = 2.80 - (1.20)^2 = 2.80 - 1.44 = \mathbf{1.36\text{ transacciones}^2}$$
  \end{alertblock}
\end{frame}

% Slide 11A: Ejercicio en Clase — Nivel 2: Operativo (Enunciado)
\begin{frame}{Ejercicio en Clase — Nivel 2: Operativo (1/2)}
  \begin{block}{Problema 3.3.4 — Licitación de Megaproyecto Minero y Retorno Ajustado por Riesgo}
    \footnotesize
    Una corporación de ingeniería evalúa un megaproyecto cuya rentabilidad neta en millones de dólares ($R$) está condicionada por escenarios globales:
    \begin{table}[h]
      \centering
      \begin{tabular}{l|ccc}
        \toprule
        \textbf{Escenario Económico} & \textbf{Probabilidad ($p_k$)} & \textbf{Rentabilidad Neta ($r_k$)} \\
        \midrule
        Expansión acelerada & $0.25$ & $+120\text{ M USD}$ \\
        Estabilidad moderada & $0.55$ & $+40\text{ M USD}$ \\
        Recesión severa & $0.20$ & $-50\text{ M USD}$ \\
        \bottomrule
      \end{tabular}
    \end{table}
    \begin{enumerate}\itemsep=0.1cm
      \item Determine el valor presente neto esperado $\E[R]$ y la desviación estándar de la rentabilidad $\sigma_R$.
      \item Si la directiva exige una métrica de seguridad financiera $U = \E[R] - 1.25\sigma_R \ge 0$, ¿debe aceptarse o rechazarse el proyecto?
    \end{enumerate}
  \end{block}
\end{frame}

% Slide 11B: Ejercicio en Clase — Nivel 2: Operativo (Resolución)
\begin{frame}{Ejercicio en Clase — Nivel 2: Operativo (2/2)}
  \fontsize{6.8pt}{8.2pt}\selectfont
  \begin{block}{1. Valor Esperado $\E[R]$ y Desviación Estándar $\sigma_R$}
    \fontsize{6.8pt}{8.2pt}\selectfont
    $$\E[R] = 120(0.25) + 40(0.55) + (-50)(0.20) = 30.0 + 22.0 - 10.0 = \mathbf{+42.0\text{ M USD}}$$
    Para obtener el riesgo $\sigma_R$, calculamos primero el momento bruto de segundo orden:
    $$\E[R^2] = (120)^2(0.25) + (40)^2(0.55) + (-50)^2(0.20) = 3600 + 880 + 500 = 4980$$
    $$\Var(R) = \E[R^2] - (\E[R])^2 = 4980 - (42.0)^2 = 3216 \implies \sigma_R = \sqrt{3216} \approx \mathbf{56.71\text{ M USD}}$$
  \end{block}

  \vspace{-0.15cm}
  \pause
  \begin{alertblock}{2. Decisión Directiva y Evaluación del Umbral de Seguridad $U$}
    \fontsize{6.8pt}{8.2pt}\selectfont
    Evaluamos el indicador de retorno ajustado por riesgo patrimonial:
    $$U = \E[R] - 1.25\sigma_R = 42.0 - 1.25(56.71) = 42.0 - 70.89 = \mathbf{-28.89\text{ M USD}}$$
    Como el resultado es negativo ($U < 0$), la volatilidad patrimonial supera la ganancia esperada. La directiva \textbf{debe rechazar formalmente} la licitación por exceso de riesgo financiero.
  \end{alertblock}
\end{frame}

% Slide 12A: Ejercicio en Clase — Nivel 3: Analítico (Enunciado)
\begin{frame}{Ejercicio en Clase — Nivel 3: Analítico (1/2)}
  \begin{block}{Problema 3.3.6 — Estandarización $Z$-Score y Momentos Invariantes}
    \footnotesize
    En una evaluación internacional de ingeniería civil con puntuación $X$, media $\mu_X$ y desviación estándar $\sigma_X > 0$, se define la variable aleatoria estandarizada ($Z$-Score):
    $$Z = \dfrac{X - \mu_X}{\sigma_X}.$$
    Demuestre formalmente y en detalle algebraico que para cualquier variable aleatoria discreta $X$, la variable estandarizada $Z$ cumple de manera rigurosa las dos identidades universales:
    \begin{enumerate}\itemsep=0.15cm
      \item Esperanza matemática nula: $\E[Z] = 0$.
      \item Varianza operacional unitaria: $\Var(Z) = 1$.
    \end{enumerate}
  \end{block}
\end{frame}

% Slide 12B: Ejercicio en Clase — Nivel 3: Analítico (Resolución)
\begin{frame}{Ejercicio en Clase — Nivel 3: Analítico (2/2)}
  \fontsize{7.5pt}{9pt}\selectfont
  Podemos reescribir la transformación como una función lineal $Z = cX + d$, donde $c = \frac{1}{\sigma_X}$ y $d = -\frac{\mu_X}{\sigma_X}$. \pause

  \vspace{-0.05cm}
  \begin{block}{1. Demostración de Esperanza Nula ($\E[Z] = 0$)}
    \fontsize{7.5pt}{9pt}\selectfont
    Aplicando linealidad universal ($\E[cX + d] = c\E[X] + d$) y recordando que $\E[X] = \mu_X$:
    $$\E[Z] = \E\left[ \dfrac{1}{\sigma_X} X - \dfrac{\mu_X}{\sigma_X} \right] = \dfrac{1}{\sigma_X}\E[X] - \dfrac{\mu_X}{\sigma_X} = \dfrac{\mu_X}{\sigma_X} - \dfrac{\mu_X}{\sigma_X} = \mathbf{0}. \quad \checkmark$$
  \end{block}

  \vspace{-0.1cm}
  \pause
  \begin{alertblock}{2. Demostración de Varianza Unitaria ($\Var(Z) = 1$)}
    \fontsize{7.5pt}{9pt}\selectfont
    Aplicando la propiedad de homogeneidad cuadrática frente a escalas ($\Var(cX + d) = c^2 \Var(X)$) y recordando que la varianza poblacional de $X$ es $\Var(X) = \sigma_X^2$:
    $$\Var(Z) = \Var\left( \dfrac{1}{\sigma_X} X - \dfrac{\mu_X}{\sigma_X} \right) = \left( \dfrac{1}{\sigma_X} \right)^2 \Var(X) = \dfrac{\Var(X)}{\sigma_X^2} = \dfrac{\sigma_X^2}{\sigma_X^2} = \mathbf{1}. \quad \blacksquare$$
  \end{alertblock}
\end{frame}

% Slide 13A: Ejercicio en Clase — Nivel 4: Desafiante (Enunciado)
\begin{frame}{Ejercicio en Clase — Nivel 4: Desafiante (1/2)}
  \begin{block}{Problema 3.3.9 — Paradoja de San Petersburgo y Regularización Truncada}
    \footnotesize
    En la Paradoja de San Petersburgo, una moneda justa se lanza repetidamente hasta obtener la primera cara en el $k$-ésimo lanzamiento ($k \ge 1$), pagándose al jugador una recompensa de $X = 2^k$ monedas.
    \begin{enumerate}\itemsep=0.12cm
      \item Demuestre que el valor esperado teórico de la recompensa diverge a infinito ($\E[X] = \infty$).
      \item Si por regulación de capital el casino impone una reserva de banca máxima de $N$ lanzamientos (si no sale cara en los primeros $N$ intentos, el juego termina pagando $M = 2^N$), obtenga una expresión cerrada para el valor esperado truncado $\E[X_N]$ y evalúela en $N = 25$ lanzamientos.
    \end{enumerate}
  \end{block}
\end{frame}

% Slide 13B: Ejercicio en Clase — Nivel 4: Desafiante (Resolución)
\begin{frame}{Ejercicio en Clase — Nivel 4: Desafiante (2/2)}
  \fontsize{7.5pt}{9pt}\selectfont
  \begin{block}{1. Divergencia Teórica de la Esperanza $\E[X]$}
    \fontsize{7.5pt}{9pt}\selectfont
    La probabilidad de obtener la primera cara en el $k$-ésimo lanzamiento es $P(X=2^k) = \left(\frac{1}{2}\right)^k$:
    $$\E[X] = \sum_{k=1}^\infty 2^k P(X = 2^k) = \sum_{k=1}^\infty 2^k \left(\dfrac{1}{2}\right)^k = \sum_{k=1}^\infty 1 = 1 + 1 + 1 + \dots = \mathbf{\infty}$$
  \end{block}

  \vspace{-0.1cm}
  \pause
  \begin{alertblock}{2. Regularización Truncada por Límite de Banca ($N=25$)}
    \fontsize{7.5pt}{9pt}\selectfont
    La variable truncada $X_N$ toma valores $2^k$ para $k \in \{1, \dots, N\}$ con probabilidad $(1/2)^k$, y el valor de cota $2^N$ con probabilidad residual $P(K > N) = (1/2)^N$:
    $$\E[X_N] = \sum_{k=1}^N 2^k \left(\dfrac{1}{2}\right)^k + 2^N \left(\dfrac{1}{2}\right)^N = \sum_{k=1}^N 1 + 1 = N + 1$$
    Para un casino con límite de $N = 25$ lanzamientos ($\approx 33.5$ millones de monedas de reserva), el precio justo de entrada es extremadamente finito y modesto:
    $$\E[X_{25}] = 25 + 1 = \mathbf{26\text{ monedas exactas}}.$$
  \end{alertblock}
\end{frame}

% Slide 14: Cierre de Sección
\begin{frame}{Cierre de Sección y Síntesis Estocástica}
  \begin{block}{Resumen de Momentos y Operadores Discretos}
    \footnotesize
    \begin{itemize}\itemsep=0.1cm
      \item \textbf{Centro de Gravedad ($\mu = \E[X]$):} Resumen de posición central obtenido ponderando el soporte por la PMF. Minimiza el error cuadrático medio poblacional.
      \item \textbf{Inercia y Riesgo ($\sigma^2 = \Var(X)$):} Medida del esparcimiento cuadrático respecto a $\mu$. La fórmula de König-Huygens ($\E[X^2] - \mu^2$) simplifica drásticamente su evaluación.
      \item \textbf{Atajos Operacionales:} El Teorema LOTUS permite evaluar transformaciones $\E[g(X)]$ sin deducir distribuciones inducidas, mientras que la estandarización $Z$ homogeniza escalas heterogéneas con $\E[Z]=0$ y $\Var(Z)=1$.
    \end{itemize}
  \end{block}
\end{frame}

% Slide 15: Perspectiva Modular
\begin{frame}{Perspectiva Modular — El Próximo Reto}
  \begin{columns}[T]
    \begin{column}{0.48\textwidth}
      \begin{block}{De Momentos Generales a Familias Específicas}
        \footnotesize
        Hasta este punto hemos desarrollado operadores de momentos aplicables a \textbf{cualquier distribución discreta arbitraria} mediante sumatorias directas sobre soportes numéricos o tabulares.
      \end{block}
    \end{column}
    \begin{column}{0.48\textwidth}
      \pause
      \begin{alertblock}{Hacia la Sección 03.04: Bernoulli y Binomial}
        \footnotesize
        En la \textbf{Sección 03.04}, especializaremos estos operadores para estudiar el experimento de Bernoulli ($X \in \{0, 1\}$) y el conteo de éxitos en $n$ ensayos independientes (Distribución Binomial), derivando analíticamente por qué $\E[X] = np$ y $\Var(X) = np(1-p)$.
      \end{alertblock}
    \end{column}
  \end{columns}
\end{frame}

\end{document}
